% Source: https://artofproblemsolving.com/community/c3h2670574_some_actually_good_proving_problems
% Stable identity: aops-topic-2670574-q12
Consider the set of all infinite sequences $x=(x_1,x_2,\ldots)$ of real numbers, and call two sequences \emph{equivalent} if they differ in only finitely many positions. This is an equivalence relation: it is clearly reflexive and symmetric, and if $x,y$ differ in a finite set $A$ of positions and $y,z$ in a finite set $B$, then $x,z$ differ only inside the finite set $A\cup B$.

\emph{Strategy.} Before the game, using the axiom of choice, the prisoners choose one representative sequence $r(C)$ from each equivalence class $C$, and all memorise this choice.

During the game the actual hats form a sequence $x$. Prisoner $P_k$ sees $x_{k+1},x_{k+2},\ldots$; the unseen entries $x_1,\ldots,x_k$ are finitely many, so all sequences agreeing with what he sees lie in the same class $C$ as $x$. Thus every prisoner knows $C$, and $P_k$ guesses the $k$-th entry of $r(C)$.

\emph{Why it works.} The representative $r(C)$ is equivalent to $x$, so $r(C)_k=x_k$ for all but finitely many $k$. Hence all but finitely many prisoners guess correctly. $\blacksquare$
